\songtitle{Réquiem por el último cassette}

\tag*{2026}\tag{indie-folk}\tag{spanish-lyrics}
\begin{notes}
Lyrics by SUNO based on a made-up lyric analysis by DeepSeek.
\end{notes}
\begin{style}
Indie folk lo-fi in Spanish, intimate central Colombian vocal pronunciation, mid-slow 6/8 groove around 78 BPM, fingerpicked nylon-string guitar with slightly imperfect timing, muted kick and brushed snare, warm upright bass, sparse melodica and distant room tone, Begin with close dry vocal and guitar; add soft tape hiss and subtle wow/flutter between sections, Build the choruses with restrained harmony vocals and fuller acoustic strums, then strip back for a fragile bridge featuring cassette button clicks and pencil-on-plastic detail, Mix warm and narrow like a worn home recording, preserving breaths, fret noise, and gentle saturation; end with the tape slowing and room ambience
\end{style}
\begin{lyrics}
\songpart{Verso 1}
Bajé por la escalera con una linterna azul,\\
el agua había dejado una raya en la pared.\\
Debajo de una cobija, dentro de un cajón de pino,\\
encontré doce cintas con tu letra y la de él.

La caja olía a humedad, a jabón y a domingo,\\
las carátulas vencidas se pegaban al pulgar.\\
Cada lado tenía un nombre, una fecha, un cumpleaños,\\
y una flecha hacia lugares que no supe conservar.

\annotation{Pre-Coro}\\
Puse la primera cinta, soplé el polvo del rodillo;\\
la canción salió torcida, pero me supo decir\\
que Julián ya no contestaba, que Marcela se casaba,\\
que el silencio también guarda lo que no quiso morir.

\annotation{Coro}\\
El ruido blanco era más honesto\\
que tu playlist de 500 canciones.\\
No prometía quedarse para siempre,\\
solo llenaba el cuarto de estaciones.\\
El ruido blanco era más honesto:\\
crujía sin venderme una versión.\\
Entre dos voces y un carrete cansado,\\
volvió la casa dentro del cajón.

\songpart{Verso 2}
Había una cinta marcada “para el viaje a Villa de Leyva”,\\
con tres canciones de lluvia y una nota sin final.\\
En otra, “no borrar”, temblaba la tinta negra;\\
la rebobiné despacio para no hacerla sangrar.

El recibo de la tienda decía: “combo familiar”,\\
como si comprar memoria fuera comprar un celular.\\
Hoy cobran por la nube, por guardar lo que olvidamos;\\
antes pagábamos monedas por poderlo rebobinar.

\annotation{Pre-Coro}\\
La grabadora tragó cinta, le metí un lápiz amarillo;\\
mis manos dieron cuerda al pequeño corazón.\\
Cada vuelta revelaba una amistad en vacaciones,\\
un beso detrás del cine, una deuda, una canción.

\annotation{Coro}\\
El ruido blanco era más honesto\\
que tu playlist de 500 canciones.\\
No prometía quedarse para siempre,\\
solo llenaba el cuarto de estaciones.\\
El ruido blanco era más honesto:\\
crujía sin venderme una versión.\\
Entre dos voces y un carrete cansado,\\
volvió la casa dentro del cajón.

\annotation{Puente}\\
Guardé las cintas secas junto al calentador,\\
pero el agua ya había escrito su propia edición.\\
La etiqueta de “nosotros” se desprendió completa;\\
quedó pegada en mi dedo como una pequeña piel.

Si mañana las vendo por internet,\\
¿quién compra el martes que te fui a buscar?\\
¿Quién pone precio al temblor de una voz\\
que ya no sabe pronunciar mi nombre igual?

\annotation{Coro}\\
El ruido blanco era más honesto\\
que tu playlist de 500 canciones.\\
No prometía quedarse para siempre,\\
solo llenaba el cuarto de estaciones.\\
El ruido blanco era más honesto:\\
crujía sin venderme una versión.\\
Entre dos voces y un carrete cansado,\\
volvió la casa dentro del cajón.

\songpart{Outro}
Apagué la grabadora, cerré la caja,\\
la dejé junto a los platos sin usar.\\
Afuera seguía lloviendo sobre el barrio;\\
adentro, una cinta aprendía a terminar.

\annotation{Instrumental}
\end{lyrics}

